\chapter{Reaction kinetics and stochastic chemistry}

Two mixtures have the same predicted equilibrium: 80 nM of duplex from
100 nM of each complementary strand. Must they produce the same signal
after ten seconds? No. Equilibrium specifies a population balance; it does
not specify how quickly that balance is approached. A molecular algorithm
has a deadline as well as a desired final state.

We will build three descriptions of one reaction. First, a concentration
equation predicts a smooth trajectory. Second, a numerical integrator
approximates that trajectory and is checked against an exact solution.
Third, an event simulator follows integer molecule counts. The third
description is not merely the first with decorative noise added.

Chapter 6 supplies mass balance and the dissociation constant. Here the
new prerequisites are derivatives as rates, basic probability, and Python
loops. Every numerical rate below is an assigned teaching parameter,
not a measured rate for a particular DNA sequence. The system is closed,
well mixed, isothermal, and restricted to distinct strands A and B and
their duplex D. Folding intermediates and competing complexes are absent.

\section{Turn a mechanism into a conservation law}

The allowed reaction channels are
\begin{equation}
 \mathrm{A+B}\xrightarrow{k_{\rm on}}\mathrm{D},
 \qquad \mathrm{D}\xrightarrow{k_{\rm off}}\mathrm{A+B}.
\end{equation}
Association consumes one copy of each single strand. Dissociation returns
both strands. Neither event copies a strand or changes its sequence.
In the species order $(A,B,D)$, the \Term{stoichiometric matrix} is
\begin{equation}
 S=\begin{pmatrix}-1&1\\-1&1\\1&-1\end{pmatrix}.
 \label{eq:stoich}
\end{equation}
Its columns describe association and dissociation. Multiplying on the left
by $(1,0,1)$ or $(0,1,1)$ gives zero, so $[A]+[D]$ and $[B]+[D]$ are conserved.
These are executable invariants, not just a bookkeeping introduction.

\Plate{01-reaction}{Two allowed molecular events. Backbone arrows mark strand polarity, not motion; dashed rungs mark pairing. Association changes counts by $(-1,-1,+1)$ and dissociation by the reverse. The plate represents species-level events, not a resolved nucleation pathway.}{fig:reaction}

Let the conserved totals be $a$ and $b$, and write $x=[D]$. Then
$[A]=a-x$, $[B]=b-x$, and $0\leq x\leq\min(a,b)$. For ideal
\Term{mass-action kinetics}, the forward and reverse fluxes are
\begin{equation}
 v_+=k_{\rm on}(a-x)(b-x),\qquad v_-=k_{\rm off}x.
 \label{eq:flux}
\end{equation}
Concentrations are mol/L, $k_{\rm on}$ is L/(mol s), and $k_{\rm off}$ is
$\mathrm{s}^{-1}$. Both fluxes therefore have units mol/(L s).
The forward channel counts possible encounters between two distinct
populations; the reverse channel acts on one existing duplex.

\Listing{fluxes}

The function checks that the proposed state respects the inventory.
An integrator that proposes $x>a$ has failed numerically; it has not
discovered a new chemical state. For the closed model,
\begin{equation}
 \frac{dx}{dt}=f(x)=k_{\rm on}(a-x)(b-x)-k_{\rm off}x.
 \label{eq:ode}
\end{equation}
At $x=0$, $f(x)\geq0$; at $x=\min(a,b)$, $f(x)\leq0$.
The exact differential equation points inward at both boundaries.
A numerical scheme is not automatically guaranteed to preserve this property.

\section{Separate equilibrium, flux, and relaxation}

Setting $f(x_*)=0$ gives the equilibrium equation from Chapter 6 with
$K_d=k_{\rm off}/k_{\rm on}$. Choose
\begin{equation}
 a=b=100\ {\rm nM},\quad k_{\rm on}=10^6\ {\rm M}^{-1}{\rm s}^{-1},
 \quad k_{\rm off}=0.005\ {\rm s}^{-1}.
\end{equation}
Then $K_d=5$ nM and $x_*=80$ nM. Initially, with no duplex,
$v_+=10$ nM/s and $v_-=0$. At equilibrium, both fluxes are
$0.4$ nM/s. The net change is zero although both reaction channels remain active.
This distinction matters when an output is sampled repeatedly or can
subsequently react with another species.

Now multiply both rate constants by ten. Their ratio is unchanged,
but the entire right-hand side of Eq.~\ref{eq:ode} is multiplied by ten.
If $x(t)$ is the original solution, the new solution is $x(10t)$
for the same initial condition. No equilibrium calculation alone can detect
this tenfold change in timing.

\Plate{02-relaxation}{Computed approach to the same equilibrium under two assigned rate pairs. The dashed curve multiplies both rates by ten. The population target is identical, but the finite-time outputs differ. The trajectories are model predictions, not measurements.}{fig:relaxation}

Close to equilibrium, write $x=x_*+\epsilon$. Linearizing gives
\begin{equation}
 \dot\epsilon\simeq-\lambda\epsilon,\qquad
 \lambda=k_{\rm on}(a+b-2x_*)+k_{\rm off},
 \quad \tau=\lambda^{-1}.
\end{equation}
For our example, $\lambda=0.045\ \mathrm{s}^{-1}$ and $\tau\approx22.22$ s.
This is a \Term{relaxation time}: a small deviation decays by a factor $e$
in one $\tau$. It is not a universal time for half the strands to bind,
nor the mean lifetime of a single duplex, which here is $1/k_{\rm off}=200$ s.
The collective return to balance depends on both channels and free-strand
availability.

\section{Build an exact trajectory before trusting an integrator}

For equal totals $a=b=C$, define the two roots
\begin{equation}
 r_\pm=\frac{2C+K_d\pm\sqrt{K_d(4C+K_d)}}2,\qquad
 r_-r_+=C^2.
\end{equation}
The physical equilibrium is $r_-$. Factorization gives
$\dot x=k_{\rm on}(x-r_-)(x-r_+)$. Separating variables,
\begin{equation}
 \frac1{r_--r_+}\ln\!\left(\frac{x-r_-}{x-r_+}\right)
   =k_{\rm on}t+\text{constant}.
\end{equation}
Applying $x(0)=0$ and defining
$q=\exp[-k_{\rm on}(r_+-r_-)t]$ yields
\begin{equation}
 x(t)=r_-\frac{1-q}{1-(r_-/r_+)q}.
 \label{eq:exact}
\end{equation}
At $t=0$ the numerator vanishes; at large $t$ the expression tends to $r_-$.
Its initial derivative is $k_{\rm on}C^2$, as required by mass action.
Those independent limits are stronger evidence than one plausible plot.

\Listing{equal_pool_exact}

The code evaluates $r_-$ as $C(C/r_+)$ to avoid subtracting similar numbers.
It evaluates $1-e^{-z}$ as $-\operatorname{expm1}(-z)$ to avoid cancellation
for small positive $z$. It deliberately supports positive rates and an
initially unbound equal mixture. Unequal pools or nonzero initial duplex
should not be silently passed through this specialized formula.

For a more general initial condition, implement the classical fourth-order
Runge--Kutta update. At each step of size $h$, evaluate four slopes:
\begin{align}
 k_1&=f(x),& k_2&=f(x+hk_1/2),\\
 k_3&=f(x+hk_2/2),& k_4&=f(x+hk_3),\\
 x_{\rm new}&=x+\frac h6(k_1+2k_2+2k_3+k_4).
\end{align}
The slopes have concentration/time units; multiplication by $h$ produces
a concentration increment. Our implementation shortens the last step to
land on the requested final time.

\Listing{rk4}

The following table compares the computed state at 20 s with
Eq.~\ref{eq:exact}. The table is generated from the reference implementation.
\begin{center}\input{\Generated/results-table.tex}\end{center}
In the small-step regime, halving $h$ reduces fourth-order global error
by approximately $2^4=16$. This is a convergence experiment, not a proof
that any requested step will be safe. The tests deliberately take a very
large step and require rejection of an invalid intermediate state.
Clipping that state into the allowed interval would conceal the failed
integration and alter the method.

Large networks may be stiff: fast and slow processes coexist, and explicit
steps may need to be extremely small for stability even when the desired
output changes slowly. This chapter supplies a small transparent oracle;
Chapter 28 will address solver choice, error tolerances, stiffness, and
parameter inference for networks.

\section{Replace concentrations with integer events}

Let $N_A$ denote Avogadro's constant and $V$ the volume in liters.
Define $\Omega=N_A V$, so a count $n$ corresponds to $n/\Omega$ mol/L.
For total strand counts $n_A,n_B$ and duplex count $d$, the channel
\Term{propensities}, measured in events/s, are
\begin{equation}
 \alpha_+(d)=\frac{k_{\rm on}}{\Omega}(n_A-d)(n_B-d),
 \qquad \alpha_-(d)=k_{\rm off}d.
 \label{eq:propensity}
\end{equation}
One association changes concentration by $1/\Omega$.
Multiplying the forward event rate by that increment recovers the
mass-action concentration drift at a fixed state. This explains the
volume factor instead of adding it as a memorized conversion.

\Listing{propensities}

For 20 copies of each strand at 100 nM,
$V=20/(N_A\,100\times10^{-9})\approx3.32\times10^{-16}$ L.
Initially $\alpha_+=2$ events/s. One duplex changes the displayed
concentration by 5 nM. A deterministic value such as 3.7 duplexes is
an aggregate prediction, never the contents of one vessel.

This is the distinct-species reaction $A+B$, not $2A$. Two identical
reactants require unordered-pair counting, $n(n-1)/2$, and a matching
rate-constant convention. Substituting $n^2$ without rederiving the
convention permits a lone molecule to react with itself.

\Plate{03-event-clock}{The next-event clock at a fixed molecular state. The total propensity determines the waiting time; normalized channel propensities determine which event occurs. Counts change only after the sampled event, then both propensities are recomputed.}{fig:clock}

\section{Implement the stochastic simulation algorithm}

At a fixed state, the probability of no reaction for duration $s$ is
$P(T>s)=e^{-\alpha_0s}$, where $\alpha_0=\alpha_++\alpha_-$.
Two independent uniform variates $u_1,u_2$ give
\begin{equation}
 \Delta t=-\frac{\ln u_1}{\alpha_0},\qquad
 \text{associate if }u_2\alpha_0<\alpha_+.
\end{equation}
Otherwise dissociate. The joint density for waiting time $s$ and channel
$j$ is $\alpha_j e^{-\alpha_0s}$: a survival factor followed by that
channel's hazard. This is the direct stochastic simulation construction
described by Gillespie \cite{gillespie}.

\Listing{ssa}

The implementation makes four boundary choices explicit. Zero total
propensity means no further events. A sampled event beyond the horizon
is not executed. Uniform zero is redrawn before taking its logarithm.
Exhausting the event budget raises an error rather than returning an
apparently complete trajectory. A local random generator keeps the seed
from changing another experiment's global random state.

The returned event list is right-continuous: after an event at $t_i$,
its new count holds until the next event. To display a trajectory at
regular observation times, select the latest event at or before each
observation. Do not linearly interpolate integer reaction events into
fictional fractional molecules.

\Plate{04-trajectory}{One seeded stochastic trajectory for 20 copies of each strand. Horizontal segments are waiting periods and vertical steps are single association or dissociation events. The final horizontal extension to 120 s is an observation horizon, not an extra event.}{fig:trajectory}

An exact event sampler means exact sampling of the specified continuous-time
Markov model, apart from random-number and floating-point limitations.
It does not mean the species list, fitted rate constants, or well-mixed
assumption are exact descriptions of real DNA.

\section{Use a distribution as a second oracle}

For positive rates, $d$ ranges over $0,\ldots,\min(n_A,n_B)$.
At stationarity, adjacent probability flows balance:
\begin{equation}
 \pi_d\alpha_+(d)=\pi_{d+1}\alpha_-(d+1),\qquad
 \pi_{d+1}=\pi_d\frac{\alpha_+(d)}{\alpha_-(d+1)}.
 \label{eq:stationary}
\end{equation}
Starting with an arbitrary weight at zero, recursively compute all
relative weights and normalize their sum. Log weights prevent intermediate
products from overflowing; subtracting their maximum does not change the
normalized probabilities.

\Listing{stationary}
\Plate{05-distribution}{Computed stationary probabilities with 20 copies of each strand, obtained independently from detailed balance. Counts 0--9 are omitted from view but retained in normalization. This describes many possible vessel states, not an empirical histogram of the single trajectory.}{fig:distribution}

An instructive exact counterexample uses only one A and one B at the
same 100 nM total concentration per species. There are two states:
$d=0$ and $d=1$. Their forward and reverse rates are 0.1 and 0.005
events/s, so
\begin{equation}
 P(d=1)=\frac{0.1}{0.1+0.005}=\frac{20}{21}\approx0.95238.
\end{equation}
The stochastic mean concentration is therefore about 95.24 nM,
not the deterministic equilibrium prediction of 80 nM. This is not
a bug to repair by forcing the simulator's mean onto the ODE.

To see why, average Eq.~\ref{eq:propensity} over the random count $D$:
\begin{equation}
 \frac{d\,E[D]}{dt}=
 \frac{k_{\rm on}}{\Omega}E[(n_A-D)(n_B-D)]-k_{\rm off}E[D].
\end{equation}
With $\mu=E[D]$,
\begin{equation}
 E[(n_A-D)(n_B-D)]=(n_A-\mu)(n_B-\mu)+\operatorname{Var}(D).
\end{equation}
The deterministic closure drops the variance term. Averaging a nonlinear
reaction rate is not the same operation as evaluating it at the average
population. Increasing volume and copy numbers at fixed concentrations
often makes relative fluctuations less important, but this needs an
appropriate limit and cannot be assumed at one or twenty molecules.

We also test the event simulator against the exact one-pair transient,
$P(D_t=1)=\frac{20}{21}(1-e^{-0.105t})$ from an unbound start.
Four thousand seeded runs are compared with that prediction using its
binomial sampling error. A fixed seed makes the regression repeatable;
the analytical oracle makes it meaningful.

\section{Know when the model boundary has been crossed}

Well mixed means spatial position is not part of the state. A diffusion
time over length $\ell$ scales as $\ell^2/D_{\rm diff}$, up to geometry
factors. If reaction and transport occur on comparable timescales,
spatial gradients, surfaces, or local depletion can matter. A molecular
association event may also conceal nucleation and zippering steps.
An effective two-state rate needs evidence that these hidden processes
can be coarse-grained for the question being asked.

The chapter's rate constants are independent assigned inputs. They are
not inferred from sequence, salt, temperature, or Chapter 6's nearest-neighbor
table. A favorable final free energy does not specify a transition barrier.
Adding a reporter, enzyme, fuel, or competing strand changes the reaction
network and may invalidate the conservation laws used here.

A computation should specify its readout event: crossing a threshold
at least once, being above threshold at a deadline, or remaining above
threshold for an interval are different observables. An equilibrium
fraction answers none of those timing questions by itself.

Run the reference and tests from the repository root:
\begin{Verbatim}[fontsize=\small]
python3 drafts/ch07/reference.py
python3 -m pytest -q tests/test_ch07_reference.py
\end{Verbatim}
The tests cover conservation, units, exact limits, fourth-order convergence,
seeded event integrity, detailed balance, and a distributional oracle.
Next, Chapter 8 introduces enzymes as substrate-specific operators.
The present inventory and timing discipline will remain necessary even
when a diagram contains a recognizable protein.

\section{Exercises with worked solutions}

\subsection*{1. Audit a reaction column}
Why is $(-1,0,+1)$ invalid for duplex formation from A and B?
\textbf{Solution:} it creates a duplex without consuming B. Its dot
product with the B-inventory vector $(0,1,1)$ is one rather than zero.

\subsection*{2. Convert a flux}
What is the initial association flux for 100 nM of each strand and
$k_{\rm on}=10^6\ \mathrm{M}^{-1}\mathrm{s}^{-1}$?
\textbf{Solution:} $10^6(10^{-7})^2=10^{-8}$ M/s, or 10 nM/s.
Using 100 as a molar concentration would be a nine-order unit error
in each concentration factor.

\subsection*{3. Explain a moving equilibrium}
At $x=80$ nM, are all reactions stopped?
\textbf{Solution:} no. Forward and reverse fluxes are both 0.4 nM/s.
Their difference is zero. Individual duplexes can still dissociate.

\subsection*{4. Rescale the clock}
Multiply both rates by three. What happens to $x(20\,\mathrm{s})$?
\textbf{Solution:} it becomes the old trajectory's value at 60 s.
The equilibrium root is unchanged and the relaxation time is divided by three.

\subsection*{5. Distinguish two lifetimes}
Why is $1/k_{\rm off}$ not the collective relaxation time?
\textbf{Solution:} it concerns an isolated duplex's dissociation hazard.
Population relaxation also includes association of available strands;
here the two times are 200 s and approximately 22.22 s.

\subsection*{6. Test numerical order}
Errors at two step sizes are $e(h)$ and $e(h/2)$. What supports fourth order?
\textbf{Solution:} a ratio near 16 across successive refinements before
roundoff dominates. One ratio at a large, unstable step is not sufficient.

\subsection*{7. Derive the volume factor}
Why divide $k_{\rm on}$ by $\Omega$ for counts?
\textbf{Solution:} the concentration flux is
$k_{\rm on}(n_A/\Omega)(n_B/\Omega)$; multiply by $\Omega$
to convert molar change into events/s. One factor of $\Omega$ remains below.

\subsection*{8. Sample a clock}
If $\alpha_0=2$ events/s and $u_1=e^{-1}$, find the waiting time.
\textbf{Solution:} $-\ln(e^{-1})/2=0.5$ s.
The next wait is redrawn after the event because the rates may change.

\subsection*{9. Respect the horizon}
The current event is at 9 s and the next sampled event at 12 s;
the output deadline is 10 s. Which count is reported?
\textbf{Solution:} the state after the 9 s event. The 12 s event must
not be applied before reporting the deadline.

\subsection*{10. Solve the one-pair equilibrium}
Use forward rate 0.1 and reverse rate 0.005 to find both probabilities.
\textbf{Solution:} $\pi_1/\pi_0=20$ and $\pi_0+\pi_1=1$,
so $\pi_0=1/21$ and $\pi_1=20/21$.

\subsection*{11. Recover the missing variance}
Expand $E[(n-D)^2]$.
\textbf{Solution:} $n^2-2n\mu+E[D^2]=(n-\mu)^2+\operatorname{Var}(D)$.
The ODE mean-field closure omits the final term.

\subsection*{12. Design a timing evaluation}
How would you evaluate a threshold-based molecular output?
\textbf{Worked rubric:} declare the species, initial inventory, volume,
rate evidence, readout rule, deadline, and acceptable error. Separate
first-passage probability from deadline occupancy. Test sensitivity to
rates and competing reactions; compare measurements with uncertainty.
A deterministic curve alone cannot establish a low-copy failure probability.
